\documentclass[10pt]{article}
\usepackage[margin=1in]{geometry}
\usepackage{amsmath,amssymb,booktabs,graphicx,hyperref}
\title{From Probabilistic Forecasts to Grid Decisions: Conformal-Scaled Wasserstein
Distributionally Robust Volt/VAR Optimization under Load and Solar-PV Uncertainty}
\author{[AUTHORS]}
\date{}

\begin{document}
\maketitle

\begin{abstract}
We investigate a forecast-to-decision framework coupling conformalized probabilistic
load/PV forecasts with Wasserstein distributionally robust Volt/VAR optimization.
[RESULT TO BE GENERATED BY FULL EXPERIMENT]. No numerical claim is inserted until the
locked full experiment and reproduction audit are complete.
\end{abstract}

\section{Introduction}
Distribution operators must choose controls before uncertain load and PV are realized. We
separate forecast accuracy, uncertainty calibration, decision quality, and decision robustness.
The investigated approach is called CS-WDRO; no novelty claim is made pending the documented
literature audit.

\section{Related Work}
Wasserstein ambiguity and CVaR/chance constraints are established in stochastic OPF,
including distribution-system formulations~\cite{guo2018stochasticopf,guo2018partone,
duan2018wassersteinopf}. Conformalized quantile regression supplies finite-sample adaptive
interval calibration~\cite{romano2019cqr}. Most importantly, conformal uncertainty sets have
already been used in robust deep-reinforcement-learning Volt/VAR control~\cite{chen2024robustvvo},
so this paper does not claim that conformal VVO is new. Recent work also analyzes the general
relationship between conformal prediction and Wasserstein DRO~\cite{long2026unified}, and
conformal forecast-to-decision methods appear in energy storage control~\cite{wu2026arbitrage}.
The investigated distinction here is the empirical behavior of forecast-specific CQR scaling
inside a convex W1 worst-case-CVaR continuous-inverter formulation with locked nonlinear AC
evaluation; the ongoing review does not establish novelty.

\section{Problem Formulation}
Define radial LinDistFlow, inverter capability, signs, objective, and pre-realization information.

\section{Probabilistic Forecasting}
Describe causal features, chronological partitions, baselines, quantile models, and raw metrics.

\section{Conformal Uncertainty Calibration}
Define finite-sample CQR, time-specific scale, and standardized two-dimensional residuals.

\section{CS-WDRO Method}
Present the W1/L1 ambiguity set and worst-case CVaR reformulation. State unrestricted-support
and affine-voltage assumptions.

\section{Experimental Design}
Predefine IEEE 33/123 studies, PV penetrations, radius selection, baselines, stress tests,
ablations, AC validation, primary endpoint, and paired block-bootstrap analysis.

\section{Results}
[TABLES AND FIGURES TO BE GENERATED BY FULL EXPERIMENT].

\section{Stress Tests and Ablations}
[RESULT TO BE GENERATED BY FULL EXPERIMENT].

\section{Discussion}
Interpret forecast accuracy, calibration, decision quality, and robustness separately. Report
negative results without post-hoc tuning.

\section{Limitations}
The study uses a linearized decision model, balanced radial feeders, low-dimensional factor
errors, mapped public time series, continuous controls, one-step decisions, simulation rather
than field deployment, and a radius selected from finite calibration data. IEEE 123 data
availability and computational scaling are additional constraints.

\section{Conclusion}
[CONCLUSION TO BE WRITTEN AFTER LOCKED ANALYSIS].

\section*{Reproducibility Statement}
Configuration snapshots, specification hashes, raw outcomes, solver states, package versions,
and generated outputs accompany every run. Synthetic quick-mode results cannot populate this paper.

\bibliographystyle{plain}
\bibliography{references}
\end{document}
